\scanpage{27}

a) On donne une partition \inserted{$\pi$} de $J$ en parties $(J_\lambda)_{\lambda\in\Lambda}$

b) Pour tout $\lambda$, on donne une \deleted{\unsure{variété}} surface orientée \inserted{connexe} à bord $U'_\lambda$, et un isom $bU'_\lambda\simeq-J_\lambda\times S^1$. Si on se borne aux $S$ ``\unsure{algébrique|algébriques}'' (compacte moins ens fini) Une telle $U'_\lambda$ sera connue quand on connait

a) le genre \inserted{$g_\lambda$} de la surface compacte sans bord obtenue en ``rajoutant les points à l'infini'' (i.e. les bouts) et en contractant chaque composante du bord en un point

b) le nb de bouts $\unsure{\beta|\nu}_\lambda$ de $U'_\lambda$.

[\uline{NB} Comme il y aura tjrs un homéom. \inserted{orienté} de $U'_\lambda$ sur lui-même \deleted{identique, \unsure{ou} isotope à l'identité} qui échange entre eux de façon prescrite les éléments du bord, on n'a pas à s'en faire pour préciser la \deleted{\unsure{nat}} correspondance entre \deleted{\unsure{pts}} \inserted{composantes} du bord et les \deleted{\unsure{éléments} de $J_\lambda$} \inserted{élém. de $J_\lambda$} \deleted{\unsure{ou}}]

En fait, en termes de $g_\lambda$, $\beta_\lambda$ on construit $U'_\lambda$ en prenant la surface \deleted{\unsure{algébr.}}\inserted{sans bord} orientée $\unsure{S}$ compacte de genre $g$, soit $S_g$, on y prend $J_\lambda$ rondelles disjointes $D_j$ ($j\in J_\lambda$), et \deleted{\unsure{un}}\inserted{une} partie \deleted{$\beta_\lambda$ points de} finie $\unsure{\Phi|P}_\lambda$ de cardinal $\beta_\lambda$ ne rencontrant pas $\bigcup_{j\in J_\lambda}D_j$, et prend $U'_\lambda=S_g-\bigcup_{j\in J_\lambda}D_j^\circ-\Phi_\lambda$
% slanted
\margin{$S$ compacte $\unsure{\Leftrightarrow}$ les $U'_\lambda$ cpcts, i.e. les $\unsure{\nu|\beta}_\lambda=0$}

Soit $S_{\alpha,\pi,g,\beta}$ la surface ainsi construite à l'aide de
\[
 \left\{\begin{aligned}
 &\alpha\in\Omega(X)\quad[\text{d'où }J=J(\alpha),\text{ à expliciter combinatoirement}]\\
 &\pi\text{ partition de }J\text{ en }(J_\lambda)_{\lambda\in\Lambda}\\
 &(g_\lambda)_{\lambda\in\Lambda}\in\deleted{\unsure{\mathbb N^{*\Lambda}}}\,\mathbb N^\Lambda\\
 &(\beta_\lambda)_{\lambda\in\unsure{J|\Lambda}}\in\deleted{\unsure{\mathbb N^{*\Lambda}}}\,\mathbb N^\Lambda
 \end{aligned}\right.
\]
On a donc un plongement canonique
\[
 X\lhook\joinrel\xrightarrow{\varphi_{\alpha,\pi,g,\beta}} S_{\alpha,\pi,g,\beta}
\]
\margin{\uline{NB} $S_{\alpha,\pi,g,\nu}$ aura $\nu=\sum_\lambda\nu_\lambda$ bouts. Il faut calculer son genre, il doit \deleted{être}\inserted{quelque chose comme} $\sum_{\lambda\in\Lambda}g_\lambda+\gamma(\alpha)$, $\gamma(\alpha)$ un entier (le genre de $U_\alpha$) \uuline{vérifier}}

On trouve donc

\uline{Th 2} Tout plongement $X\longrightarrow S$ \inserted{($S$ surface orientée \uline{algébrique})} est isomorphe à un $\varphi_{\alpha,\pi,g,\beta}$ pour $\alpha,\pi,g,\beta$ convenables, uniquement déterminé
% archive sheet number: 26
